\documentclass[../IC.tex]{subfiles}

\begin{document}
	\chapter{Conjugação Topológica}
	
	O objetivo deste capítulo é relacionar o mapa shift discutido anteriormente com o mapa logístico $F_\mu$, quando $\mu > 2 + \sqrt{5}$. Lembre-se de que com exceção dos pontos no conjunto de Cantor $\Lambda$, todos os outros pontos em $\mathbb{R}$ tendem a $-\infty$ sob a iteração de $F_\mu$. Para completar a descrição da dinâmica de $F_\mu$, devemos então compreender a restrição de $F_\mu$ para $\Lambda$.
	
	\section{O Mapa Itinerário}
	
	O mapa itinerário será um mapa do conjunto de Cantor para o espaço de sequências $\Sigma$. Para definir esse mapa, voltamos para $\Lambda \subset I_0 \cup I_1$. Se $x \in \Lambda$, então todos pontos na órbita de $x$ também estão em $\Lambda$, e com isso estão em um desses dois intervalos. Podemos, desse modo, obter uma ideia aproximada do comportamento da órbita observando em qual desses intervalos caem as várias iteradas de $x$. Assim, estabelecemos a seguinte definição.
	
	\begin{definicao}
		O itinerário de $x$ é a sequência $S(x) = s_0s_1s_2...$ onde $s_j = 0$ se $F_\mu^j(x) \in I_0$ e $s_j = 1$ se $F_\mu^j(x) \in I_1$.
	\end{definicao}
	
	Assim, o itinerário de $x$ é uma sequência infinita de 0s e 1s. Ou seja, $S(x)$ é um ponto no espaço de sequências $\Sigma_2$. Vemos $S$ como um mapa que nos leva de $\Lambda$ a $\Sigma_2$.
	
	\begin{teorema}
		Se $\mu > 2+\sqrt{5}$, então o mapa itinerário $S: \Lambda \rightarrow \Sigma_2$ é um homeomorfismo.
	\end{teorema}
	
	\begin{prova}
		conteúdo...
	\end{prova}
	
	\section{Conjugação}
	
	O teorema acima mostra que, como conjuntos, $\Lambda$ e $\Sigma_2$ são a mesma coisa. A relação existente entre os mapas $F_\mu$ e $\sigma$ é que eles são topologicamente conjugados.
	
	\begin{definicao}
		Seja $f: A \rightarrow A$ e $g: B \rightarrow B$ dois mapas. $f$ e $g$ são topologicamente conjugados se existe um homeomorfismo $h: A \rightarrow B$ tal que $h \circ f = g \circ h$. A equação $h \circ f = g \circ h$ é dita equação de conjugação e o homeomorfismo $h$ é dito conjugação topológica.
	\end{definicao}
	
	Mapas que são topologicamente conjugados são completamente equivalentes em termos de dinâmica. Pontos periódicos e fixos, órbitas e densidade são preservados na transição. Desse modo, estudar a dinâmica em $g$ é a mesma coisa que estudar a dinâmica em $f$.
\end{document} 